\section{Candidate propositions}

\subsection{Candidate 1: finite-window reachability comparison}
\mstatus{CONDITIONAL} \istatus{PARTIAL} \nstatus{NOT CLAIMED} The following
is the strongest surviving candidate.

\begin{proposition}[Memory advantage on a stable target window]
\label{prop:window}
Suppose the true hard region attached to the declared discovery event is
constant, $\HARD_j=\HARD$, for a window of $W$ comparison opportunities and
each opportunity has the same deterministic compute cost
$c>0$. On every no-hit history in that window, suppose a memory policy $M$
has conditional hit hazard at least $p_M$, and a comparator $G$ has
conditional hit hazard at most $p_G$, with $p_M>p_G$. Then, for
$1\leq n\leq W$,
\[
 \Prb(\tau_M\leq n)\geq 1-(1-p_M)^n
 > 1-(1-p_G)^n\geq \Prb(\tau_G\leq n),
\]
and, for $2\leq n\leq W$,
\[
 \Ex[c\min(\tau_M,n)]
 \leq c\sum_{j=0}^{n-1}(1-p_M)^j
 < c\sum_{j=0}^{n-1}(1-p_G)^j
 \leq \Ex[c\min(\tau_G,n)].
\]
For a generic policy that actually samples a parent, one sufficient
derivation of $p_M$ is the following: on each no-hit history,
with conditional probability at least $r$ a useful retained parent is
available; given that event it is selected with probability at least
$\gamma$; conditional on that selection, its mutation/evaluation reaches the
true current hard region with probability at least $\alpha$. Then
$p_M\geq r\gamma\alpha$. No independence among these events is assumed;
the bounds are conditional in the stated order.

That random-selection factorization is not the S6-AE schedule. For the
verified deterministic Cartesian batch, every retained parent receives all
11 operators. If $r$ lower-bounds availability of a useful parent and
$\alpha_{\rm batch}$ lower-bounds, under the explicitly declared RNG
ensemble, the conditional probability that its scheduled operator batch
generates the declared hit, then the appropriate sufficient bound is
$p_M\geq r\alpha_{\rm batch}$. With roots and all seeds conditioned to their
frozen values, this is a statement about a realized batch, not an estimated
random parent/operator-selection probability.
\end{proposition}

\paragraph{Importance.} It turns a structural-memory claim into a finite-budget,
equal-cost comparison that explicitly accounts for parent access and mutation
reachability. It also makes clear that survival alone is not a sufficient
mechanism.

\paragraph{Assumptions and realism.} A piecewise-stable target window and
equal-cost opportunities are useful first approximations. In S6-AE the BQ
checkpoint is fixed during mutation/acquisition in cycle $k$ and changes only
after $T_k$ is frozen and labeled, but comparison windows spanning adaptations
can still move $\HARD_k$. HGB40, archive, exact-label, and BQ-training costs
also differ from a single mutation attempt. The deterministic 10,000-by-11
schedule and $P=T$ update are verified; the history-uniform
$r\alpha_{\rm batch}$ bound, scientific threshold, and matched comparator
cost/hazard remain \istatus{OPEN}.

\paragraph{Proof approach.} Induction on the no-hit survival probability,
then the identity expressing expected truncated hitting time as the sum of
survival probabilities. It requires no conditional independence and no
Markov approximation. Proof details are in \texttt{proof\_attempts.tex}.

\paragraph{Counterexample and HAC distinction.} If memory's conditional
hazard is lower because its retained parents are stale, or if a gradient
proposal has larger hazard per compute unit, the premise fails. The exact
examples in \texttt{counterexamples.tex} show both. Thus this is a conditional
criterion that can distinguish BLADE from fixed generation or HAC only after
their hazards and costs are measured; it does not establish that BLADE wins.

\subsection{Candidate 2: inherited memory never hurts}
\mstatus{DISPROVED} \istatus{NOT APPLICABLE} \nstatus{NOT CLAIMED} The
proposed blanket claim is that retaining historical
parents cannot increase first-hitting cost relative to a fixed generator.
This is false even with two structural states and one target change. The
stale-memory example in Section~\ref{sec:counterexamples} has expected cost
$100$ for inherited memory versus $2$ for a fixed uniform generator. The
failure is not a technical gap in Candidate 1: the required hazard dominance
is reversed.

\paragraph{Importance and proof approach.} A universal no-harm theorem would
have supported the research objective directly, but the exact counterexample
rules it out. No proof is possible without additional target-drift,
forgetting, or reachability assumptions. These extra assumptions would need
to be stated and independently justified; they cannot be inferred from the
word ``memory.''

\paragraph{Assumptions and realism.} The universal claim would require memory never to reduce true-hardness reachability per unit cost for any target shift or parent-selection history. That is not guaranteed by inherited populations and is contradicted by the stale-memory example; realism for BLADE-NCO is low absent explicit target-aware forgetting and verified hazard bounds.

\paragraph{Strongest proof approach.} An exact finite-state counterexample proves the universal claim false. A positive restricted result would need additional drift and reachability premises, which are instead stated explicitly in Candidate 1.

\paragraph{HAC distinction.} Nothing in this false claim compares BLADE to
HAC. HAC can outperform a poorly aligned retained population, as the path
example in Section~\ref{sec:counterexamples} demonstrates.

\subsection{Candidate 3: diversity always improves discovery}
\mstatus{DISPROVED} \istatus{NOT APPLICABLE} \nstatus{NOT CLAIMED} The
proposed blanket claim is that preserving a more
diverse population always reduces first-hitting time. In a three-state model,
a diverse parent selector spends half of each equal-cost proposal on an
irrelevant niche and has expected hitting cost $2$, while a focused selector
hits in cost $1$. The exact construction is in Section~\ref{sec:counterexamples}.

\paragraph{Importance and assumptions.} Diversity can preserve useful future
routes, but its benefit depends on niche relevance, selection allocation,
target drift, and cost. A monotonicity proof from diversity cardinality alone
cannot work. This counterclaim concerns a universal guarantee, not a claim
that Elite-QD1000 behaves like the toy selector.

\paragraph{Realism and strongest proof approach.} No selection-allocation or niche-relevance condition is supplied for Elite-QD1000, so the premise is not verified for BLADE-NCO. The exact three-state construction is a complete disproof of the universal claim; a positive diversity theorem would require explicit target-access conditions. This does not rank BLADE against HAC.

\paragraph{HAC distinction.} It provides no general ranking against HAC.
With equal compute, a focused or gradient-guided policy may spend more
proposals on the current reachable hard region; the actual ranking depends
on transition probabilities and costs, neither of which is specified here.
